\chapter{Conclusões}
\label{chap:conclusoes}

O trabalho concretiza um protótipo que liga preparação de dados, previsão
de resultados, avaliação estatística de jogadores e cinco iniciais,
planeamento de rotações, substituições e objetivos de vários jogos.
A divisão pelos seis elementos articula as três componentes de AAUTIA
com as três componentes de PLANTIA através de dados, modelos e verificações
comuns.

\section{Principais resultados}

Na previsão de resultados, a regressão logística selecionada atingiu
67,2\% de exatidão e perda logarítmica de 0,603 no teste. Superou a
frequência histórica de vitória em casa e apresentou qualidade
probabilística ligeiramente superior ao Elo, embora este tenha obtido
maior exatidão. A avaliação individual por Ridge melhorou pouco o MAE
face à média histórica. Para os cinco iniciais, o RMSE de 15,503 pontos
foi inferior aos 16,407 da referência, mantendo uma capacidade explicativa
limitada e uma interpretação associativa.

A programação por restrições produziu uma rotação de 240 minutos-jogador
com verificação independente e prova de optimalidade para a formulação
do cenário. O MDP preservou os orçamentos de minutos, mas não melhorou a
utilidade das políticas fixa e gulosa no exemplo avaliado. Este resultado
não sustenta uma vantagem empírica da adaptação nesse cenário.
O algoritmo genético, aplicado às políticas do mesmo MDP, igualou a solução
exata nas cinco execuções. Na sensibilidade sem penalização de fadiga,
ambos superaram ligeiramente as políticas fixa e gulosa. Contudo, o espaço
de apenas oito cromossomas não permite concluir superioridade ou melhor
escalabilidade da pesquisa genética.

A integração selecionou rotações para três jogos sob limites acumulados
e reportou 77,1\% de probabilidade de atingir duas vitórias, assumindo
independência. Essa probabilidade pertence ao modelo de resultados;
não representa um ganho provocado pelo plano. A principal realização da
integração é compatibilizar previsões, restrições e políticas sequenciais
num procedimento reproduzível.

\section{Limitações e trabalho futuro}

As fontes e os modelos não fornecem uma avaliação causal de jogadores ou
decisões. Disponibilidade, posições, calendário histórico e minutos da
fonte exigem validação adicional. Os cinco iniciais não substituem a
reconstrução dos intervalos efetivos em campo. A fadiga, os contextos e
os pesos de utilidade permanecem hipóteses de modelação.

Os próximos passos incluem validar a codificação dos minutos e os horários
dos relatórios de lesões, reconstruir alinhamentos com atribuição fiável
de tempo e pontos e repetir a avaliação em diferentes temporadas.
No planeamento, importa estudar cenários com maior liberdade de ação,
calendários consecutivos e indisponibilidades, quantificar sensibilidade
às previsões e calibrar transições com dados adequados. Só uma avaliação
adicional poderá sustentar afirmações sobre melhoria de decisões ou de
resultados desportivos.
